\section{Dynamic Ramanujan mean and scalar-tail removal}
\label{sec:dynamic-mean}

Two scalar issues must be settled before the conductor argument.  The
principal character of the enlarged packet has to remain harmless, and
the centered coefficient beyond \(S\) has to be replaced by its actual
divisor sum without a formal deletion.

\subsection{The mean packet for
\texorpdfstring{\(u,v\le S\)}{u,v <= S}}

Let \(\cZ_1=\sum_{m,n}|z_{m,n}|\) be the absolute row-pair mass in one
separated packet.  The application has
\begin{equation}
 \cZ_1\ll Q^2X^\eps.
 \label{eq:row-l1}
\end{equation}
For each compatible pair \(u,v\le S\), let
\(\mathfrak R_{u,v}\) denote its row-pair mass after the primitive and
row-pair masks.  If \(q=[u,v]\), the unit-fiber mean is
\begin{equation}
 \overline Z_{q,S}
 =\frac1{\varphi(q)}
 \sum_{\substack{u,v\le S\\\operatorname{lcm}(u,v)=q}}
 \vartheta_{R,S}(u)\vartheta_{R,S}(v)\mathfrak R_{u,v}.
 \label{eq:dynamic-fiber-mean}
\end{equation}
This is exact and uses no row equidistribution.

For \(M\ge1\), define
\begin{equation}
 \Delta_M(F;J)
 =\frac JM\sum_{r\ne0}c_M(r)
   \widehat F\left(\frac{rJ}{M}\right).
 \label{eq:Delta}
\end{equation}
Poisson summation over reduced classes gives
\begin{equation}
 \Delta_M(F;J)
 =\sum_{(n,M)=1}F(n/J)
  -\frac{\varphi(M)}M J\widehat F(0),
 \qquad |\Delta_M(F;J)|\ll_F\tau(M).
 \label{eq:Delta-bound}
\end{equation}
The mean scalar packet is exactly
\begin{equation}
 \cL_{q,S}^{\rm mean}=\overline Z_{q,S}\Delta_{Hq}(F;J).
 \label{eq:mean-physical}
\end{equation}

\begin{lemma}[Dynamic lcm cost]
\label{lem:dynamic-lcm-cost}
For \(S\le X^{O(1)}\) and every \(\eps>0\),
\begin{equation}
 \boxed{
 \sum_{u,v\le S}
 \frac{|\vartheta_{R,S}(u)\vartheta_{R,S}(v)|
       \tau(H[u,v])}{\varphi([u,v])}
 \ll_{\eps,H}X^\eps.}
 \label{eq:dynamic-lcm-cost}
\end{equation}
\end{lemma}

\begin{proof}
By \eqref{eq:b-soft}, each coefficient is \(O(X^{\eps/4})\) after
shrinking the auxiliary epsilon.  Also, uniformly for
\(q\le S^2\le X^{O(1)}\),
\[
 \frac{\tau(Hq)}{\varphi(q)}\ll_{\eps,H}\frac{X^{\eps/4}}q.
\]
Writing \(u=ga,v=gb\), then dropping \((a,b)=1\), gives
\begin{equation}
 \sum_{u,v\le S}\frac1{[u,v]}
 \le\sum_{g\le S}\frac1g
 \left(\sum_{a\le S/g}\frac1a\right)^2
 \ll\log^3(2S).
 \label{eq:reciprocal-lcm}
\end{equation}
Absorb the logarithms and coefficient factors into \(X^\eps\).
\end{proof}

\begin{theorem}[Dynamic mean closure]
\label{thm:dynamic-mean}
For any collection of denominators from the selected packet
\(u,v\le S\),
\begin{equation}
 \boxed{
 \sum_q|\cL_{q,S}^{\rm mean}|
 \ll_{\eps,F,H}\cZ_1X^\eps.}
 \label{eq:dynamic-mean-bound}
\end{equation}
Under \eqref{eq:row-l1}, this is \(O(Q^2X^\eps)\), and therefore has
a fixed-power saving relative to \(XQ=Q^2J\) whenever \(J\) has a
fixed positive power of \(X\).
\end{theorem}

\begin{proof}
Insert \eqref{eq:dynamic-fiber-mean} into
\eqref{eq:mean-physical}, apply \eqref{eq:Delta-bound}, and bound the
masked compatible row mass absolutely by \(\cZ_1\).  Interchanging the
row and divisor sums leaves exactly \eqref{eq:dynamic-lcm-cost}.
\end{proof}

This proof is a genuine extension of the TPC-21 mean theorem: that
theorem stated only the supports \(u\le R\).  The extension works
because \(b_R(u)=a(u)\) for \(u>R\), not because the principal block is
discarded.

\subsection{Uniform removal of the ultra-long scalar tail}

Put
\begin{equation}
 A_T(m)=\sum_{\substack{u\le T\\(u,mH)=1}}\frac{a(u)}u,
 \qquad \rho_{P,H}(m)=\frac{mH}{\varphi(mH)}.
 \label{eq:AT}
\end{equation}
The uniform M\"obius calibration theorem of TPC-20, based only on the
classical zero-free region for \(\zeta(s)\), says that for every fixed
\(A,C>0\), uniformly for \(mH\le T^C\),
\begin{equation}
 A_T(m)=\rho_{P,H}(m)
 +O_{A,C}\left(\rho_{P,H}(m)(\log T)^{-A}\right).
 \label{eq:uniform-calibration}
\end{equation}
In the present fixed-power ranges, one fixed \(C\) covers both
\(T=S\) and \(T=U\asymp X\).

\begin{proposition}[The centered ultra-long tail is its divisor tail]
\label{prop:scalar-tail}
Define
\begin{equation}
 \tau_{S,U}(m)=
 \sum_{\substack{S<u\le U\\(u,mH)=1}}\frac{a(u)}u.
 \label{eq:scalar-tail}
\end{equation}
Then, for every \(A>0\),
\begin{equation}
 \tau_{S,U}(m)
 \ll_{A,h}\rho_{P,H}(m)(\log X)^{-A}
 \ll_{A,h,\eps}X^\eps(\log X)^{-A}.
 \label{eq:scalar-tail-small}
\end{equation}
Consequently
\begin{equation}
 Z_m^{>S}(j)=\cT_{m,S}(j)-\tau_{S,U}(m),
 \qquad
 \cT_{m,S}(j)=\sum_{\substack{u\mid n_m(j)\\u>S}}a(u).
 \label{eq:centered-to-divisor-tail}
\end{equation}
Replacing every \(Z_m^{>S}\) by \(\cT_{m,S}\) in a separated two-row
packet changes it by \(O_B(XQ(\log X)^{-B})\) for every prescribed
\(B>0\).
\end{proposition}

\begin{proof}
Equation \eqref{eq:scalar-tail} is \(A_U(m)-A_S(m)\), so
\eqref{eq:scalar-tail-small} follows by applying
\eqref{eq:uniform-calibration} at both endpoints.  Divisors of
\(n_m(j)\) are automatically coprime to \(mH\), which proves
\eqref{eq:centered-to-divisor-tail}.

For the packet error, the elementary bounds
\begin{equation}
 |Z_m^{\le S}(j)|+|\cT_{m,S}(j)|
 \ll_\eps X^\eps
 \label{eq:pointwise-signals}
\end{equation}
follow from the divisor bound, \eqref{eq:b-soft}, and
\(\sum_{u\le U}|b_R(u)|/u\ll(\log X)^2\).  Expand the product after
inserting \eqref{eq:centered-to-divisor-tail}, use \(O(Q^2J)\) row--orbit
tuples and \(Q^2J\asymp XQ\), then choose \(A\) in terms of \(B\).
\end{proof}

\begin{remark}[Calibration is not cancellation of the parity core]
\label{rem:calibration-not-parity}
The proposition removes only the deterministic densities \(a(u)/u\).
It does not estimate the signed divisor tail \(\cT_{m,S}\).  In
particular, its \(u=n_m(j)\) term remains present whenever the target is
squarefree.
\end{remark}
